% Number_Sense_Place_Value.tex
\documentclass[12pt]{article}
% Engine-aware Unicode / Korean (Hangul) support.
% Recommended: compile with XeLaTeX or LuaLaTeX. A pdfLaTeX fallback is provided.
\usepackage{iftex}
\ifPDFTeX
  % pdfLaTeX branch (legacy): use CJKutf8
  \usepackage[utf8]{inputenc}
  \usepackage[T1]{fontenc}
  \usepackage{CJKutf8}
  % Usage for pdfLaTeX: \kr{...} wraps Korean text using a CJK family code.
  \newenvironment{foreigntext}[1]{\begin{CJK}{UTF8}{#1}}{\end{CJK}}
  \newcommand{\kr}[1]{\begin{foreigntext}{mj}#1\end{foreigntext}}
\else
  % XeLaTeX / LuaLaTeX branch (recommended)
  \usepackage{fontspec}
  \usepackage{xeCJK}
  % Set fonts (change to fonts installed on your system)
  \setmainfont{TeX Gyre Termes}
  \setsansfont{TeX Gyre Heros}
  \setCJKmainfont{Noto Serif CJK KR}
  % Simple wrapper: under Xe/Lua this is identity
  \newcommand{\kr}[1]{#1}
\fi

% Usual math / layout packages
\usepackage{amsmath,amssymb}
\usepackage{booktabs}
% \usepackage{array}



\usepackage{geometry}
\geometry{margin=1in}
\title{Basic Number Concepts: Place Value, Decimals \\ \& Fractions}
\author{}
\date{}

\begin{document}
\maketitle
\maketitle

\section*{Overview}
This document converts the provided notes into a clean LaTeX source. It covers:
\begin{itemize}
  \item Place value (whole numbers and decimals),
  \item Decimals and fractions (terminology and conversions),
  \item Comparing and ordering numbers (vocabulary and symbols),
  \item Prime and composite numbers (prime factorization and divisibility tests),
  \item Greatest common factor (GCF) and least common multiple (LCM) with the ``\texttt{\kr{ㄴ}}'' method.
\end{itemize}

\section{Place Value}
\subsection{Essential whole-number place values}
The important whole-number place values to learn (from largest to smallest) are:
\begin{center}
\begin{tabular}{@{}lll@{}}
\toprule
English & Korean & Value (power of 10) \\
\midrule
billions & \kr{십억} (or \kr{억}, depending on grouping) & $10^{9}$ \\
hundred millions & \kr{백억} & $10^{8}$ \\
ten millions & \kr{십억} & $10^{7}$ \\
millions & \kr{백만} & $10^{6}$ \\
hundred thousands & \kr{십만} & $10^{5}$ \\
ten thousands & \kr{만} (\kr{만} = $10^{4}$ in Korean system) & $10^{4}$ \\
thousands & \kr{천} & $10^{3}$ \\
hundreds & \kr{백} & $10^{2}$ \\
tens & \kr{십} & $10^{1}$ \\
ones & \kr{일} & $10^{0}$ \\
\bottomrule
\end{tabular}
\end{center}

\textbf{Note on Korean grouping:} In Korean the grouping differs from English 
after the thousands place. Korean uses units of \kr{만} ($10^{4}$), \kr{억} ($10^{8}$), \kr{조} ($10^{12}$), etc. The user notes: ``\kr{한국어와} \kr{비슷하게} \kr{앞자리에} 10\kr{배만} \kr{붙이면} \kr{됨}: \kr{일} \kr{십} \kr{백} \kr{천} \kr{만} \kr{십만} \kr{백만} \kr{천만} \kr{억} \kr{십억} \kr{백억}'' — this is a mnemonic showing successive 10× steps and how the Korean names appear.

\subsection{Decimal place values}
Decimal places (to the right of the decimal point) are named by adding ``-ths'':
\begin{itemize}
  \item tenths (\emph{\kr{할}} in Korean) $= 10^{-1}$,
  \item hundredths (\emph{\kr{푼}}) $= 10^{-2}$,
  \item thousandths (\emph{\kr{리}}) $= 10^{-3}$,
  \item ten-thousandths $= 10^{-4}$, and so on.
\end{itemize}

The notes recommend memorizing tenths, hundredths, and thousandths especially.

\section{Decimals and Fractions}
\subsection{Terminology}
\begin{description}
  \item[Numerator] (\kr{분자}): the top part of a fraction.
  \item[Denominator] (\kr{분모}): the bottom part of a fraction.
\end{description}

If remembering the words is difficult, the user suggests the mnemonic ``AND'' to recall the order numerator/denominator (``and'' = n/d order).

\subsection{Types of fractions}
\begin{itemize}
  \item \textbf{Improper fraction (\kr{가분수})}: numerator $\ge$ denominator (e.g. $7/4$).
  \item \textbf{Mixed number (\kr{대분수})}: an integer plus a proper fraction (e.g. $1\tfrac{3}{4}$).
\end{itemize}
The suggestion: understanding the property is usually enough to remember the terms.

\subsection{Types of decimals}
\begin{itemize}
  \item \textbf{Repeating / recurring decimal (\kr{순환소수})}: digits repeat periodically, often notated with a bar (e.g. $1.\overline{3}=1.333\ldots$).
  \item \textbf{Finite decimal (\kr{유한소수})}: a decimal with a finite number of nonzero digits after the point (e.g. $0.25$).
  \item \textbf{Infinite decimal (\kr{무한소수})}: non-terminating and non-repeating decimals; these are generally not expressible as fractions.
\end{itemize}

\subsubsection{Detecting repeating decimals in notation}
A repeating decimal is often indicated by placing a dot or bar over the repeating digits. Examples:
\begin{itemize}
  \item $1.\dot{3} = 1.333\ldots = \dfrac{4}{3}$.
  \item If the digits $3475$ repeat in $1.34753475\ldots$, it may be written with dots or a vinculum over the repeating block.
\end{itemize}

\subsubsection{Converting a repeating block to a fraction}
If a repeating block has $k$ digits and starts immediately after the decimal point, the repeating portion equals an integer over $10^{k}-1$ (a string of $k$ nines). For example, if $x=0.\overline{3984}$ then
\[x=\dfrac{3984}{9999}.\]
If the repeating block starts after some non-repeating digits, handle by subtracting the non-repeating portion or using standard algebraic techniques (noted above as the usual method).

\section{Comparing and Ordering Numbers}
Common comparison symbols and English vocabulary:
\begin{center}
\begin{tabular}{@{}ll@{}}
$>$ & greater than (more than, larger than, higher than, above, exceeds, over) \\
$\ge$ & greater than or equal to (at least, minimum, no less than) \\
$<$ & less than (fewer than, smaller than, lower than, under, below) \\
$\le$ & less than or equal to (at most, maximum, no more than) \\
\end{tabular}
\end{center}

The notes suggest focusing on the most test-relevant phrases (e.g. ``at least/at most'') and to read each problem carefully to interpret the intended meaning.

\section{Prime and Composite Numbers}
\subsection{Prime factorization}
Prime factorization is the process of breaking a positive integer into a product of prime numbers (its smallest multiplicative building blocks).

\textbf{Basic procedure:}
\begin{enumerate}
  \item Check small prime divisors (2, 3, 5, 7, \ldots) using divisibility tests.
  \item Divide repeatedly by primes until the remaining factor is $1$.
\end{enumerate}

\subsection{Common divisibility tests (useful heuristics)}
\begin{itemize}
  \item \textbf{Divisible by 2:} last digit is even (0,2,4,6,8).
  \item \textbf{Divisible by 3:} sum of digits is divisible by 3.
  \item \textbf{Divisible by 4:} last two digits form a number divisible by 4 (or end with 00).
  \item \textbf{Divisible by 5:} last digit is 0 or 5.
  \item \textbf{Divisible by 7:} a quick test is: for a three-digit grouping $abc$, check whether $ab-2c$ is divisible by 7 (this is optional / rarely needed).
\end{itemize}

\subsection{Example: Factorize 60}
\[60\div 2 = 30,\quad 30\div 2 = 15,\quad 15\div 3 = 5,\quad 5\div 5 = 1.\]
Hence the prime factorization is
\[60=2\times 2\times 3\times 5=2^{2}\cdot 3\cdot 5.\]

\section{Greatest Common Factor (GCF) \& Least Common Multiple (LCM)}
The notes describe the Korean schoolbook ``\texttt{\kr{ㄴ}}'' method for 
finding GCF and LCM of two or more numbers. The idea is to successively 
divide the numbers by common divisors and record the divisors.

\subsection{Procedure (``\texttt{\kr{ㄴ}}'' method) and example}
Find the GCF and LCM of 12 and 18.

Start by arranging the numbers side-by-side and dividing by a common factor repeatedly (here 2 then 3):
\[\begin{array}{r|rr}
   6 & 12 & 18 \\
   \cline{2-3}
     & 2 & 3 \\
\end{array}\]

In this representation the last row (to the right/below the dividing line) shows the quotients 
after dividing by the common factors; when the remaining numbers are prime (or no larger 
common factor remains) the process ends.

\textbf{GCF:} the product of the common divisors used in the vertical division process. In this case the common divisors are $2$ and $3$, so
\[\mathrm{GCF}(12,18)=2\times 3=6.\]

\textbf{LCM:} the product of the GCF and any remaining quotients (equivalently, 
multiply all numbers that appear on the right and bottom of the \texttt{\kr{ㄴ}} 
chart). Using the example above,
\[\mathrm{LCM}(12,18)=6\times 2\times 3=36.\]
(Another standard method: compute via prime factorizations and take the highest power of each prime that appears.)

\section*{Closing Notes}
This LaTeX source faithfully converts the supplied learning notes into structured sections and examples. If the user prefers a different document class (\texttt{article} vs \texttt{report}), additional Korean font packages (for guaranteed correct Korean rendering), or a compiled PDF, indicate the preference and a downloadable file can be prepared.

\end{document}
